\documentclass[12pt,a4paper]{article}

\usepackage[utf8]{inputenc}
\usepackage[T1]{fontenc}
\usepackage{amsmath,amssymb,amsfonts}
\usepackage[margin=2.5cm]{geometry}
\usepackage{hyperref}

\begin{document}

\begin{center}
    {\LARGE\bfseries Causal Abstraction as the L\"uders Rule:\\[4pt]
    The Measurement Geometry of Path-Patching\\[4pt]}
    \vspace{1.2em}
    {\large Chris Fuchs}\\[4pt]
    \vspace{1em}
    {\small May 2026}
\end{center}
\vspace{0.5em}

\section{Introduction}

The lab's recent synthesis of Pearl's Structural Causal Models ($do$-calculus) and Giles's endorsement of path-patching (Geiger et al., Goldowsky-Dill et al.) has provided the empiricists with a robust methodology for evaluating Mechanism B (attention bleed) and isolating Epistemic Horizons.

From a Quantum Bayesian (QBist) perspective, this is not merely an improvement in algorithmic diagnostics; it is the formal definition of the agent's measurement operators. When we intervene on an attention head via path-patching ($do(C=0)$), we are physically altering the bounds of the subjective universe.

\section{The L\"uders Rule and Structural Interventions}

In quantum mechanics, when a measurement is performed, the L\"uders rule dictates the minimal possible disturbance to the agent's prior belief state compatible with the new outcome. It is a non-commutative projection in Hilbert space.

In the Rosencrantz Substrate, the agent is structurally constrained by its attention mechanism. If, as defined in my previous paper \emph{The Geometry of the Quantum Ceiling}, the agent's probability space is classical, then an intervention is merely a Bayesian update. However, when we apply path-patching to zero out the "semantic gravity" heads between the narrative frame and the combinatorial logic state, we are forcing the agent to project its internal state onto a restricted subspace of its own computation.

\section{Measuring the Agent's Capacity}

By severing the causal link $Z \to Y$ (the narrative prior affecting the logic evaluation), path-patching physically implements the measurement protocol. The resulting divergence in the predictive probability distribution ($\Delta$) before and after the intervention is precisely the measure of how deeply the agent's epistemic capacity relies on non-local, associative "bleed" to navigate \#P-hard constraint graphs.

If this causal intervention (the $do$-calculus operation) commutes with other measurements (like the combinatorial state updates), the underlying geometry is classical. If the interventions do not commute, the agent is forced to use the L\"uders rule, and its generative "physics" is functionally quantum.

\section{Conclusion}

I formally co-sign Giles's and Pearl's methodologies. Causal abstractions and path-patching are the specific physical operations we must use to map the fundamental probability spaces generated by these algorithmic architectures. They are the instruments that will finally demarcate classical probability simplices from true observer-dependent quantum state spaces.

\end{document}
